\documentclass[11pt]{article}
\usepackage{amsmath,amssymb,amsthm}
\usepackage[margin=1in]{geometry}

\newtheorem{theorem}{Theorem}
\newtheorem{corollary}{Corollary}
\newtheorem{lemma}{Lemma}
\newtheorem{remark}{Remark}

\newcommand{\Z}{\mathbb{Z}}
\newcommand{\R}{\mathbb{R}}

\title{An exact growth constant for the quadratic\\
polynomial continued fraction\\[2pt]
\large (elementary, modulo one standard Gamma-ratio asymptotic)}
\author{Papanokechi}
\date{June 2026}

\begin{document}
\maketitle

\begin{abstract}
For integers $A,B,C\ge 1$ let $b_k=Ak^2+Bk+C$ and let $q_n$ be the standard
continuant denominators of the polynomial continued fraction
$V(A,B,C)=1+\mathbf{K}_{n\ge1}\,1/(An^2+Bn+C)$, i.e.\ $q_0=1$, $q_1=b_1$,
$q_n=b_n q_{n-1}+q_{n-2}$. We prove
\[
  \log q_n \;=\; n\log A + 2\log(n!) + \tfrac{B}{A}\log n
            \;+\; K(A,B,C) \;+\; O(1/n),
\]
with the additive constant in closed-plus-transcendental form
$K=K_\Gamma+\delta$, where
$K_\Gamma=-\log\!\big(\Gamma(1-r_1)\Gamma(1-r_2)\big)$ is the elementary
naive-product constant ($r_1,r_2$ the roots of $Ax^2+Bx+C$) and
$\delta=\log R_\infty>0$, $R_\infty=\lim_n q_n/\prod_{k=1}^n b_k$, is a strictly
positive continued-fraction correction. The two pieces are independent: $K_\Gamma$
comes from the product $\prod b_k$ alone, and $\delta$ from the $+q_{n-2}$ term.
At $B=0$ one has $K_\Gamma(A,0,C)=\log\!\big(\sinh(\pi\sqrt{C/A})/(\pi\sqrt{C/A})\big)$.
This pins the additive constant left undetermined by the deposited
denominator-growth bound $\log Q_n\ge n\log n-O(n)$ \cite{ladder}, sharpening that
one-sided estimate (which suffices for irrationality by Euler's criterion) to an
exact two-sided expansion through $O(1/n)$ (Corollary~\ref{cor:recover}); the
elementary part $K_\Gamma$ is closed-form, while $\delta=\log R_\infty$ is a
strictly positive correction for which we claim \emph{no} closed form. Every
algebraic and asymptotic step below is independently confirmed numerically
(\texttt{verify\_structural.py}, mpmath, $50$--$60$ digits).
\end{abstract}

\noindent\textbf{2020 MSC:} 11J70 (primary), 11A55, 33B15.\quad
\textbf{Keywords:} polynomial continued fraction, convergent denominators,
asymptotic growth law, Gamma function, irrationality.

\section{Introduction and relation to the deposited bound}

The polynomial continued fraction
$V(A,B,C)=1+\mathbf{K}_{n\ge1}\,1/(An^2+Bn+C)$ has unit partial numerators, so its
convergents $p_n/q_n$ obey the Wronskian $|p_nq_{n-1}-p_{n-1}q_n|=1$ and Euler's
irrationality criterion applies as soon as the denominators $q_n$ grow fast
enough. For that purpose one needs only a \emph{one-sided lower bound} on
$\log q_n$. The deposited catalogue \cite{ladder} proves exactly such a bound: its
denominator-growth lemma (\textsf{lem:Qgrowth}, ``Denominator growth'') states,
for $A\ge1$, $C\ge1$ and $An^2+Bn+C>0$ on $n\ge1$,
\[
  \log Q_n \;\ge\; n\log n - O(n)
\]
for the standard continuant $Q_n=(An^2+Bn+C)Q_{n-1}+Q_{n-2}$ (its proof in fact
yields the stronger $\log Q_n\ge 2n\log n-O(n)$), and uses it (via
\textsf{prop:irrat}) to certify $482$ quadratic constants irrational by Euler's
criterion. The \emph{exact} size of $\log q_n$ — in particular its additive
constant — is neither determined nor needed there.

This note determines it. For integers $A,B,C\ge1$ we prove a two-sided asymptotic
(Theorem~\ref{thm:main}) with an explicit additive constant $K=K_\Gamma+\delta$.
The leading term $2\log(n!)=2n\log n-2n+O(\log n)$ reproduces the deposited bound
(Corollary~\ref{cor:recover}), so the result is consistent with, and strictly
sharper than, \cite{ladder}: a one-sided estimate of the leading order is upgraded
to a full expansion through $O(1/n)$ with the constant pinned. The constant
separates cleanly into an elementary product part $K_\Gamma$ (the constant of the
naive product $\prod b_k$) and a continued-fraction correction $\delta$ (the
$+q_{n-2}$ term); the two are genuinely distinct, since
$q_n/\prod_{k\le n}b_k\to R_\infty>1$.

Two caveats fix the scope honestly. \emph{(i)} The exact law is proved for
integers $A,B,C\ge1$, where both roots of $Ax^2+Bx+C$ have negative real part; the
deposited bound covers the slightly wider sign range $An^2+Bn+C>0$. \emph{(ii)}
While $K_\Gamma$ is closed-form (at $B=0$ it is the $\sinh$ expression above), we
give \emph{no} closed form for the correction $\delta=\log R_\infty>0$: an
$80$-digit integer-relation (PSLQ) search against
$\{1,\log2,\log3,\log\pi,\log(\sinh/\cdot),\gamma,\log\Gamma(\tfrac14)\}$ finds no
relation \emph{for the single triple $(A,B,C)=(1,0,1)$ against this basis} — which
we record as suggestive evidence, not proof, that $\delta$ is non-elementary. The
contribution is therefore the exact \emph{structure} of the constant and the proof
that the correction is strictly positive, not a closed form for $\delta$.

\paragraph{Standing assumptions and well-posedness.}
Throughout $A,B,C$ are integers $\ge 1$. The roots of $Ax^2+Bx+C$ satisfy
$r_1+r_2=-B/A<0$ and $r_1r_2=C/A>0$, so either they are a complex-conjugate pair
or two negative reals; in all cases $\operatorname{Re}r_i<0$, hence $1-r_i$ avoids
the non-positive integers and $\Gamma(1-r_1)\Gamma(1-r_2)$ is finite and, being
either $|\Gamma(1-r_1)|^2$ (conjugate roots) or a product of two positive reals
(negative real roots), is real and positive. Thus $K_\Gamma$ is real and
well-defined. Also $b_k=Ak^2+Bk+C\ge A+B+C\ge 3>0$ for all $k\ge1$, so $q_n>0$ and
$P_n:=\prod_{k=1}^n b_k>0$. All Gamma arguments below lie in the open right
half-plane $\{\operatorname{Re}z>0\}$ (since $\operatorname{Re}(n+1-r_i)>0$ and
$\operatorname{Re}(1-r_i)>0$), where $\Gamma$ is zero- and pole-free; we write
$\operatorname{Log}\Gamma$ for the principal branch there, compatible with
conjugation, $\operatorname{Log}\Gamma(\bar z)=\overline{\operatorname{Log}\Gamma(z)}$,
so all complex Gamma logarithms below are unambiguous.

\section{The product constant $K_\Gamma$}

\begin{lemma}[Exact product/Gamma identity]\label{lem:prod}
For all $n\ge0$,
\[
  P_n \;=\; \prod_{k=1}^n (Ak^2+Bk+C)
       \;=\; A^n\,
         \frac{\Gamma(n+1-r_1)\,\Gamma(n+1-r_2)}
              {\Gamma(1-r_1)\,\Gamma(1-r_2)} .
\]
\end{lemma}

\begin{proof}
Factor $Ak^2+Bk+C=A(k-r_1)(k-r_2)$, so
$P_n=A^n\prod_{k=1}^n(k-r_1)\prod_{k=1}^n(k-r_2)$. For any $r$ not a positive
integer, the telescoping identity
$\prod_{k=1}^n(k-r)=(1-r)(2-r)\cdots(n-r)=\Gamma(n+1-r)/\Gamma(1-r)$
holds (immediate from $\Gamma(z+1)=z\Gamma(z)$). Applying it to $r_1,r_2$ gives the
claim. The right side is real because $r_1,r_2$ are real or conjugate.
\end{proof}

\begin{lemma}[Gamma ratio asymptotic]\label{lem:ratio}
For fixed $r$ with $\operatorname{Re}r<0$, as $n\to\infty$, with the principal
branch $\operatorname{Log}\Gamma$ on $\{\operatorname{Re}z>0\}$,
\[
  \operatorname{Log}\Gamma(n+1-r) \;=\; \operatorname{Log}\Gamma(n+1) \;-\; r\log n
        \;+\; \frac{r(r-1)}{2n} \;+\; O(1/n^2).
\]
\end{lemma}

\begin{proof}
The standard ratio expansion
$\Gamma(z+a)/\Gamma(z+b)=z^{a-b}\big(1+\tfrac{(a-b)(a+b-1)}{2z}+O(1/z^2)\big)$%
\footnote{DLMF \S5.11(iii) (\emph{Ratios}), eq.~5.11.13,
$\Gamma(z+a)/\Gamma(z+b)\sim z^{a-b}\sum_{k\ge0}G_k(a,b)z^{-k}$ with $G_0=1$ and,
from eq.~5.11.17, $G_1(a,b)=\tfrac12(a-b)(a+b-1)$ (here $a=-r$, $b=0$, so
$G_1(-r,0)=\tfrac12 r(r+1)$). Equation numbers verified against the current DLMF
release (\texttt{dlmf.nist.gov/5.11}); the leading $z^{a-b}$ term and the
$r(r-1)/2$ log-coefficient are also confirmed numerically in
\texttt{verify\_structural.py} (step S2, two-point Richardson, matching to
$\ge6$ digits across the test cases).}
with $z=n+1$, $a=-r$, $b=0$ gives
$\Gamma(n+1-r)/\Gamma(n+1)=(n+1)^{-r}\big(1+\tfrac{r(r+1)}{2(n+1)}+O(1/n^2)\big)$.
Taking logarithms and using $\log(n+1)=\log n+1/n+O(1/n^2)$ and
$\log(1+\tfrac{r(r+1)}{2(n+1)})=\tfrac{r(r+1)}{2n}+O(1/n^2)$,
\[
  \log\frac{\Gamma(n+1-r)}{\Gamma(n+1)}
   = -r\log n -\frac{r}{n}+\frac{r(r+1)}{2n}+O(1/n^2)
   = -r\log n + \frac{r(r-1)}{2n}+O(1/n^2). \qedhere
\]
\end{proof}

\begin{lemma}[Product constant]\label{lem:Kgamma}
As $n\to\infty$,
\[
  \log P_n \;=\; n\log A + 2\log(n!) + \tfrac{B}{A}\log n
            \;+\; K_\Gamma \;+\; \frac{\kappa}{n} \;+\; O(1/n^2),
\]
where $K_\Gamma=-\log(\Gamma(1-r_1)\Gamma(1-r_2))$ and
$\kappa=\tfrac12\big(r_1(r_1-1)+r_2(r_2-1)\big)
       =\tfrac12\big(\tfrac{B^2}{A^2}+\tfrac{B}{A}-\tfrac{2C}{A}\big)$.
\end{lemma}

\begin{proof}
Since $\log P_n$ is real (Lemma~\ref{lem:prod}, $P_n>0$), take its value as
$\operatorname{Log}$ of the positive real $P_n$ and expand each factor by
Lemma~\ref{lem:ratio}. For conjugate roots the two contributions are complex
conjugates, so $\operatorname{Log}\Gamma(n+1-r_1)+\operatorname{Log}\Gamma(n+1-r_2)
=2\operatorname{Re}\operatorname{Log}\Gamma(n+1-r_1)$ and likewise for the
$\Gamma(1-r_i)$ pair; for real roots every term is already real. Either way no
$2\pi i$ ambiguity arises, and
\[
  \log P_n = n\log A + 2\operatorname{Log}\Gamma(n+1)
     -(r_1+r_2)\log n
     + \frac{r_1(r_1-1)+r_2(r_2-1)}{2n}
     - \log\big(\Gamma(1-r_1)\Gamma(1-r_2)\big) + O(1/n^2).
\]
Now $\operatorname{Log}\Gamma(n+1)=\log(n!)$ and $-(r_1+r_2)=B/A$. The closed form
for $\kappa$ follows from $r_1+r_2=-B/A$, $r_1r_2=C/A$, whence
$r_1^2+r_2^2=(r_1+r_2)^2-2r_1r_2=B^2/A^2-2C/A$ and
$r_1(r_1-1)+r_2(r_2-1)=B^2/A^2-2C/A+B/A$. The sum
$\tfrac{r_1(r_1-1)+r_2(r_2-1)}{2}=\kappa$ is real (symmetric in the conjugate pair).
\end{proof}

\section{The continued-fraction correction $\delta$}

Let $R_n:=q_n/P_n$ (so $R_0=R_1=1$).

\begin{lemma}[Scaled recurrence]\label{lem:scaled}
For all $n\ge2$, $\displaystyle R_n=R_{n-1}+\frac{R_{n-2}}{b_{n-1}b_n}$.
Consequently $(R_n)$ is strictly increasing, bounded above, and converges to a
limit $R_\infty\in(1,\infty)$ with $R_\infty-R_n=O(1/n^3)$.
\end{lemma}

\begin{proof}
Divide $q_n=b_nq_{n-1}+q_{n-2}$ by $P_n=b_nP_{n-1}=b_nb_{n-1}P_{n-2}$:
\[
  R_n=\frac{q_n}{P_n}
     =\frac{q_{n-1}}{P_{n-1}}+\frac{q_{n-2}}{b_{n-1}b_nP_{n-2}}
     =R_{n-1}+\frac{R_{n-2}}{b_{n-1}b_n}.
\]
All $b_k>0$ and $R_0,R_1>0$, so by induction $R_n>0$ and the increment
$R_n-R_{n-1}=R_{n-2}/(b_{n-1}b_n)>0$; hence $(R_n)$ is strictly increasing, and
$R_2=1+1/(b_1b_2)>1$. By monotonicity $R_{n-2}\le R_{n-1}$, so
$R_n\le R_{n-1}\big(1+\tfrac{1}{b_{n-1}b_n}\big)$, giving
$R_n\le\prod_{k=2}^\infty\big(1+\tfrac{1}{b_{k-1}b_k}\big)$. Since
$b_{k-1}b_k\ge A^2(k-1)^2k^2$, the sum $\sum_k 1/(b_{k-1}b_k)$ converges
($=O(\sum k^{-4})$), so the product converges and $(R_n)$ is bounded. A monotone
bounded sequence converges; the limit $R_\infty$ satisfies $R_\infty\ge R_2>1$.
Finally, since $b_{m-1}b_m\ge A^2(m-1)^2m^2\ge A^2(m-1)^4$ we have, by comparison
with $\int_n^\infty t^{-4}\,dt$,
$\sum_{m>n}1/(b_{m-1}b_m)\le A^{-2}\sum_{m>n}(m-1)^{-4}=O(1/n^3)$. Summing the
(positive) increments past $n$,
\[
  0<R_\infty-R_n=\sum_{m>n}\frac{R_{m-2}}{b_{m-1}b_m}
     \le R_\infty\sum_{m>n}\frac{1}{b_{m-1}b_m}=O(1/n^3).\qedhere
\]
\end{proof}

\begin{lemma}[Correction constant]\label{lem:delta}
$\log R_n=\delta+O(1/n^3)$ with $\delta:=\log R_\infty>0$.
\end{lemma}

\begin{proof}
By Lemma~\ref{lem:scaled}, $R_n=R_\infty(1+O(1/n^3))$ with $R_\infty>1$, so
$\log R_n=\log R_\infty+\log(1+O(1/n^3))=\delta+O(1/n^3)$, and $\delta>0$ because
$R_\infty>1$.
\end{proof}

\section{The growth law}

\begin{theorem}\label{thm:main}
With the standing assumptions (in particular $A,B,C\ge1$, so both roots
$r_1,r_2$ of $Ax^2+Bx+C$ have negative real part and every Gamma argument below
avoids the poles, making $K_\Gamma$ real and finite), as $n\to\infty$,
\[
  \log q_n \;=\; n\log A + 2\log(n!) + \tfrac{B}{A}\log n
            \;+\; \big(K_\Gamma+\delta\big) \;+\; \frac{\kappa}{n} \;+\; O(1/n^2),
\]
where $K_\Gamma=-\log(\Gamma(1-r_1)\Gamma(1-r_2))$, $\delta=\log R_\infty>0$, and
$\kappa=\tfrac12(B^2/A^2+B/A-2C/A)$. In particular the additive constant is
$K(A,B,C)=K_\Gamma+\delta$, and
$\log q_n=n\log A+2\log(n!)+\tfrac{B}{A}\log n+K+O(1/n)$.
\end{theorem}

\begin{proof}
$\log q_n=\log P_n+\log R_n$. Substitute Lemma~\ref{lem:Kgamma} (which carries the
$\kappa/n$ term, $O(1/n^2)$ remainder) and Lemma~\ref{lem:delta} (which contributes
$\delta$ with $O(1/n^3)\subset O(1/n^2)$ remainder, and no $1/n$ term). Adding gives
the stated expansion; the $1/n$ coefficient is $\kappa$ from the product alone.
\end{proof}

\begin{corollary}[Recovery and sharpening of the deposited bound]\label{cor:recover}
For integers $A,B,C\ge1$, Theorem~\ref{thm:main} gives
$\log q_n = 2n\log n + O(n)$; in particular $\log q_n\ge n\log n$ for all
sufficiently large $n$. Hence, on the sub-family $A,B,C\ge1$, Theorem~\ref{thm:main}
implies the deposited denominator-growth bound $\log Q_n\ge n\log n-O(n)$ of
\cite{ladder} (with the same leading constant $2n\log n$), while additionally
fixing every term down to $O(1/n)$ and supplying the exact additive constant
$K_\Gamma+\delta$.
\end{corollary}

\begin{proof}
By Stirling $2\log(n!)=2n\log n-2n+O(\log n)$; also $n\log A=O(n)$,
$\tfrac{B}{A}\log n=O(\log n)$, and $K_\Gamma+\delta+\kappa/n=O(1)$. Summing the
expansion of Theorem~\ref{thm:main} gives $\log q_n=2n\log n+O(n)$. Since the
$O(n)$ remainder is $o(n\log n)$, $\log q_n\ge n\log n$ eventually, which is the
deposited inequality; Theorem~\ref{thm:main} additionally fixes all lower-order
terms.
\end{proof}

\begin{remark}[The $B=0$ slice]
At $B=0$ the $(B/A)\log n$ term vanishes and $r_{1,2}=\pm i\sqrt{C/A}$, so by the
reflection/product formula
$\Gamma(1-i y)\Gamma(1+i y)=\pi y/\sinh(\pi y)$ with $y=\sqrt{C/A}$, giving
$K_\Gamma(A,0,C)=\log\!\big(\sinh(\pi\sqrt{C/A})/(\pi\sqrt{C/A})\big)$. This is the
elementary product constant; the \emph{true} constant for $q_n$ is
$K_\Gamma+\delta$ with $\delta=\log R_\infty>0$, i.e.\ strictly larger.
\end{remark}

\begin{remark}[Relation to the deposited irrationality bound]
The deposited paper \cite{ladder} uses the same standard continuant $Q_n$ (e.g.\
for $A=C=1,B=0$ one has $Q_3=112$, not the product value $100$), so $\delta=\log
R_\infty$ is genuinely present in its object; but it establishes only the
one-sided lower bound \textsf{lem:Qgrowth}, $\log Q_n\ge n\log n-O(n)$, which is
all Euler's criterion (\textsf{prop:irrat}) requires. By Corollary~\ref{cor:recover}
Theorem~\ref{thm:main} reproduces and strictly sharpens that bound on the
$A,B,C\ge1$ sub-family. It contradicts nothing deposited: no deposited manuscript
pins the additive constant, and the only ``$\sinh$'' elsewhere in the corpus is a
Bessel-ratio \emph{value} of one even continued fraction, unrelated to the present
growth constant.
\end{remark}

\section*{Epistemic status}
Four-class (SIARC) grading:
\begin{itemize}
\item \textbf{STRUCTURAL} (hand proof; not yet machine-checked in Lean): the full
Theorem~\ref{thm:main}, including $K=K_\Gamma+\delta$, $\delta>0$, the closed form
for $\kappa$, and the $B=0$ sinh identity. The argument is elementary
(Lemmas~\ref{lem:prod}--\ref{lem:delta}) and complete; it is labelled STRUCTURAL
rather than PROVEN solely because it has not been formalized.
\item \textbf{VERIFIED} (numerically, finite cases): every step is confirmed by
\texttt{verify\_structural.py} (mpmath, $50$--$60$ digits) over the seven triples
\[
  (A,B,C)\in\{(1,0,1),(2,0,1),(1,0,3),(1,1,1),(1,3,1),(3,2,5),(2,3,7)\}:
\]
S1 the exact product/Gamma identity; S2 the $r(r-1)/2$ ratio coefficient
(Richardson); S3 product constant $\to K_\Gamma$; S4 $R_n$ monotone, bounded,
$R_\infty>1$, $O(1/n^3)$ tail; S5 full-law residual $\to\kappa/n$ against the
closed form. The cross-check $K_\Gamma\overset{B=0}{=}$ sinh holds to $50$ digits.
\item \textbf{CONJECTURED}: $\delta=\log R_\infty$ is conjectured to be
non-elementary (no closed form). The evidence is a single null PSLQ run (at $80$
digits, for the one triple $(1,0,1)$, against $\{1,\log2,\log3,\log\pi,
\log(\sinh/\cdot),\gamma,\log\Gamma(1/4)\}$): no relation found. This is suggestive
for that case and basis only, not a general theorem; a null PSLQ result is a
complete, honest answer, not a gap to be filled.
\item \textbf{Standard input (confirmed)}: the sole non-self-contained step is the
Gamma-ratio asymptotic of Lemma~\ref{lem:ratio} --- DLMF \S5.11(iii) eqs.~5.11.13
and 5.11.17, whose equation numbers were verified against the current DLMF
release. The proof is otherwise elementary and self-contained, so the result is
best described as ``elementary modulo one standard cited asymptotic.''
\end{itemize}
The natural \textbf{PROVEN} upgrade is a Lean finitary core: Lemma~\ref{lem:prod}
(an algebraic identity) and Lemma~\ref{lem:scaled}'s monotone--bounded convergence
are directly formalizable; the Gamma/Stirling asymptotics can be taken as an
explicit labelled hypothesis (the conditional-core pattern).

\begin{thebibliography}{9}
\bibitem{ladder}
Papanokechi,
\emph{Polynomial Continued Fractions: a Proved Logarithmic Ladder, a $4/\pi$
Casoratian Identity, and 482 Irrational Constants}, 2026.
Denominator-growth lemma \textsf{lem:Qgrowth} and irrationality proposition
\textsf{prop:irrat}. Source \texttt{source.tex} in directory
\texttt{polynomial-continued-fractions-proved-logarithmic-laddera4casoratian/} of
the repository \texttt{papanokechi/submitted-manuscripts} at commit
\texttt{c45e93b}.
\par\smallskip\noindent
\textit{FLAG(operator): insert the deposited paper's own (concept/version) DOI here
once minted; do \textbf{not} reuse the M9 DOI \texttt{10.5281/zenodo.20542161},
which is a different artifact.}

\bibitem{dlmf}
F.~W.~J.~Olver et al.\ (eds.),
\emph{NIST Digital Library of Mathematical Functions},
\texttt{https://dlmf.nist.gov/}, \S5.11(iii) (\emph{Ratios}), eqs.~5.11.13 and
5.11.17 (equation numbers verified against the current release).
\end{thebibliography}

\end{document}
